\section*{Data Management and Sharing Plan}

\subsection*{Products and Data Types}

The project will produce source code, schemas, tests, documentation,
containers and Spack recipes, software bills of materials (SBOMs), benchmark
results, tutorial notebooks, synthetic datasets, Frozen Legacies catalog and
dataset manifests, observation/flight-line GeoJSON, configuration/provenance
manifests, and de-identified service-usage aggregates. Science demonstrations
will use model inputs/outputs and publicly available observational products.
Individual experiment outputs may be large and remain in user-owned project or
HPC storage; the project will preserve the smaller inputs, configurations,
diagnostics, and metadata required to understand and reproduce published
results. The project will not collect sensitive human-subject research data.

\subsection*{Formats, Metadata, and Quality Control}

Source and documentation will use Python, shell, YAML, JSON/JSON Schema,
Markdown/MyST, Jupyter notebooks, and LaTeX. Scientific arrays will use
NetCDF/HDF5 or Zarr with CF-style variable names, dimensions, coordinates,
units, fill values, and attributes where applicable. Tabular aggregates will
use CSV with a data dictionary. Each release will record semantic version,
source commit, dependency locks, image digest, SBOM, supported platform,
license, test results, and known limitations. Each qualified historical-data
collection will record source, steward, license, attribution,
spatial/temporal coverage, coordinate reference system, file inventory and
checksums, transformations, uncertainty, and expert-reviewed quality-control
status. Each reproducible run bundle will record application,
dataset-adapter/model version, observation source and checksum, configuration,
random seeds, resource request, backend, timestamps, environment reference,
output checksums, and diagnostic definitions.

Automated unit, schema, regression, numerical, restart, end-to-end, and
portability tests will gate releases. Synthetic test data will be generated
from documented scripts with fixed seeds. Scientific acceptance tolerances and
reference outputs will be versioned and reviewed by the science lead.

\subsection*{Access, Licensing, and Reuse}

Code, schemas, documentation, tests, and nonrestricted tutorial data will be
released no later than the associated public release through the
ICESEE-project GitHub organization. Tagged releases will be archived in Zenodo
or an equivalent long-term repository with DOI and CITATION.cff metadata. The
project will harmonize currently inconsistent BSD/MIT metadata in Year 1. The
intended default for original CryoStack/ICESEE code is the permissive BSD
2-Clause License, subject to Georgia Tech approval, because it supports
academic and commercial reuse while preserving attribution. Documentation and
tutorial prose will use CC BY 4.0 where institutionally permitted. Third-party
models, data, MATLAB components, and dependencies retain their own licenses;
the distribution will include a machine-readable dependency/license manifest
and will not redistribute restricted inputs.

Observation bundles and Frozen Legacies records will cite the original DOI/URL
or archive identifier and preserve source license, required attribution,
collection lineage, and documented transformations. Data too large or restricted for redistribution will
be represented by a fetch script, checksum, metadata record, and access
instructions. Publications will link to the exact release and archived
artifacts needed to support their conclusions.

\subsection*{Usage, Privacy, and Security}

Operational records will include a random experiment identifier, application,
adapter/backend versions, coarse resource request, status, timestamps, and
errors needed for reliability and metrics. Credentials, private keys,
passwords, and access tokens will not be stored in the research dataset or
public metrics. Account email and optional affiliation will be separated from
run telemetry with access limited to authorized project personnel. Public
reports will use aggregate counts, suppress small demographic cells, exclude
automated tests, and document definitions. Raw identifiable operational logs
will be retained only as long as required for service security and annual
report verification, then deleted according to a written retention schedule.

\subsection*{Preservation and Responsibilities}

GitHub provides active version control, issue history, and review records;
Zenodo provides archival releases and persistent identifiers. Annual metric
extracts, analysis scripts, data dictionaries, benchmark reports, and final
governance/operations documents will be deposited by the end of each reporting
year. Georgia Tech systems will back up the production database and deployment
configuration; restoration will be tested annually. The PI is responsible for
software releases, provenance, security-sensitive access, and repository
preservation. The Co-PI is responsible for scientific-data selection,
scientific validation criteria, and citation completeness. Adapter maintainers
are responsible for model-specific metadata and test fixtures. The team will
review this plan annually and document material changes in project reports.
